% TOP_PROBLEM: {"problemId": "problem.polynomial-time-unknot-recognition", "canonicalTitle": "Polynomial-time unknot recognition", "releaseRank": 317, "primaryDomain": "geometry_topology", "primaryDomainLabel": "Geometry & topology", "displayStatus": "open", "releaseStatus": "open"}
% !TeX program = lualatex
% Definition and sources only; this file is not an LLM solution attempt.
\documentclass[11pt]{article}
\usepackage[margin=25mm]{geometry}
\usepackage{fontspec}
\setmainfont{Latin Modern Roman}
\usepackage{amsmath,amssymb,mathtools}
\usepackage{unicode-math}
\setmathfont{Latin Modern Math}
\usepackage{xurl,hyperref}
\hypersetup{colorlinks=true,urlcolor=blue}
\setcounter{secnumdepth}{0}
\setlength{\parindent}{0pt}
\setlength{\parskip}{5pt}
\setlength{\emergencystretch}{3em}
\begin{document}
\section*{\texorpdfstring{Polynomial-time unknot recognition}{Polynomial-time unknot recognition}}\label{317}
\subsection{Definitions and mathematical statement}
A knot diagram is a finite planar projection of an embedded circle, with transverse double points and an over/under choice at each crossing. Its encoding records the planar combinatorics and crossing information; let $n$ be the bit length of that encoding. The unknot is the ambient-isotopy class of a round circle in $S^3$. The task is to decide whether there exist a deterministic algorithm and constants $C,k$ such that, on every valid diagram $D$, the algorithm correctly decides
\[
 D\text{ represents the unknot}
\]
in at most $C(n+1)^k$ steps. The complexity concerns the finite diagram input, not real coordinates with unit-cost exact operations. Certificates, special families of diagrams, and practical heuristics do not by themselves supply a polynomial-time decision algorithm for all diagrams.
\par\smallskip\textit{Formulation and concept references:} \href{https://arxiv.org/abs/1211.1079}{[S1]}.

\subsection{Short English statement}
Is there a deterministic algorithm that recognizes the unknot from any knot diagram in time polynomial in the diagram's encoding length?
\subsection{Sources}
\begin{itemize}
\item[S1]Burton and Ozlen, arXiv:1211.1079, 2012; revised 2014.
\url{https://arxiv.org/abs/1211.1079}.
\textit{Formulation source.}
\end{itemize}
\end{document}
